\scanpage{21}

pond respectively to the $\tau$ which \unsure{with} $1$ \unsure{generate} \unsure{the one} lattice of algebraic integers of $\mathbb{Q}(i)$ (i.e. $\mathbb{Z}+i\mathbb{Z}$), and the lattice of integers of $\mathbb{Q}(j)$, made up of numbers of the form $aj+b\bar{j}$ ($j=\dfrac{-1+i\sqrt{3}}{2}=\exp\dfrac{2i\pi}{3}$); the \unsure{corresponding} ramification is $2$ and $3$.

We have a canonical (analytic) isomorphism \inserted{\unsure{between}}
\[
\tag{29}
D/\mathrm{GP}(1,\mathbb{Z})^{+}\simeq\mathbb{P}^{1\,\mathrm{an}}\setminus\{0\}=Y,
\]
\unsure{characterized} by the condition that the critical orbit of ramification $3$ \unsure{correspond} to the point $\infty$, and that of ramification $2$ to the point $1$. We thus find an isomorphism
(putting $\mathcal{G}=\mathrm{GP}(1,\mathbb{Z})$, $G=\mathrm{GP}(1,\mathbb{Z})^{+}=\mathrm{Sl}(2,\mathbb{Z})'$)

\begin{samepage}
(30)
\diagram{
\begin{tikzcd}[ampersand replacement=\&,column sep=large]
\pi_1(D,\mathcal{G},x)
  \arrow[r,"\sim"] \arrow[d,equal]
\& \pi_1^{\unsure{\pm|*}}(Y,\underline{d};y) \\
\mathcal{G} \&
\end{tikzcd}
}
\end{samepage}
\margin{$(x\in D\setminus\unsure{D^{!}|D'})$}

which I would like to make explicit, by identifying \deleted{the} $\tau'_0,\tau'_1,\tau'_\infty$ as \unsure{elements} of $\mathrm{Gl}^{\unsure{\pm'|\pm}}(2,\mathbb{Z})'$.

\margin{\unsure{0}}
\margin{20}

\scanpage{22}
\mspage{299}

\uline{39} \deleted{\unsure{Properties}}\inserted{The} congruence conditions
for the outer operation of $\Gamma_{\mathbb Q}$ on $\widehat{\Pi}_{0,3}$.

Let us consider the diagram

\[
\tag{1}
\diagram{\begin{tikzcd}[ampersand replacement=\&,column sep=large,row sep=small]
1 \arrow[r] \& \widehat{\mathcal T}_{1,1}^{+} \arrow[r] \arrow[d,phantom,"\simeq" description,sloped] \& N_{1,1} \arrow[r] \arrow[dd] \& \Gamma_{\mathbb Q} \arrow[r] \arrow[dd,"\chi"] \& 1 \\
\& \mathrm{Sl}(2,\mathbb Z)^{\wedge} \arrow[d] \& \& \& \\
1 \arrow[r] \& \mathrm{Sl}(2,\widehat{\mathbb Z}) \arrow[r] \& \mathrm{Gl}(2,\widehat{\mathbb Z}) \arrow[r,"\det"] \& \widehat{\mathbb Z}^{*} \arrow[r] \& 1
\end{tikzcd}}
\]

It shows that the outer operation of
$\Gamma_{\mathbb Q}$ on $\widehat{\mathcal T}_{1,1}^{+}$
\inserted{$\simeq \mathrm{Sl}(2,\mathbb Z)^{\wedge}$}
passes \unsure{to the|to the} quotient as an
outer action on $\mathrm{Sl}(2,\widehat{\mathbb Z})$, and this
outer action, moreover, can be made
entirely explicit with the help of the
cyclotomic character
\[
\tag{2}
\chi : \Gamma_{\mathbb Q}\longrightarrow \widehat{\mathbb Z}^{*}.
\]

To make this last point precise, let us note
that the second row of (1) gives
a homomorphism
\[
\tag{3}
\widehat{\mathbb Z}^{*}\longrightarrow
\mathrm{Autext}(\mathrm{Sl}(2,\widehat{\mathbb Z})),
\]
and the canonical homomorphism
\[
\tag{4}
\Gamma_{\mathbb Q}\longrightarrow
\deleted{\unsure{\mathrm{Sl}(2,}}
\mathrm{Autext}(\mathrm{Sl}(2,\widehat{\mathbb Z}))
\]
\margin{21}

\scanpage{23}

is none other than the composite of (2) and (3).
Moreover, the homomorphism (4) factors
as
\[
\tag{5}
\widehat{\mathbb Z}^{*}/\widehat{\mathbb Z}^{*2}
\longrightarrow \mathrm{Autext}(\mathrm{Sl}(2,\widehat{\mathbb Z}))
\]

(\unsure{comes} from the fact that the \unsure{elements|elements} of
$\mathrm{Gl}(2,\widehat{\mathbb Z})$ which are in the center
$\widehat{\mathbb Z}^{*}$ operate trivially
\deleted{\unsure{from the point}}\inserted{on} $\mathrm{Sl}(2,\widehat{\mathbb Z})$,
and that the composite
\[
\widehat{\mathbb Z}^{*}\xrightarrow{\text{(center)}}
\mathrm{Gl}(2,\widehat{\mathbb Z})\xrightarrow{\det}
\unsure{\widehat{\mathbb Z}|\mathbb Z}^{\unsure{*|2}}
\]
is none other than $\lambda\longmapsto\lambda^{\unsure{2}}$). More
precisely, we see that for the
\unsure{automorphism} of $\mathrm{Sl}(2,\widehat{\mathbb Z})$ induced by a
$u\in\mathrm{Gl}(\unsure{1|2},\widehat{\mathbb Z})$ to be inner,
\unsure{i.d.|i.e.} of the form
$\mathrm{int}\,\unsure{(u_0)|u_0}$ with $u_0\in\mathrm{Sl}(2,\widehat{\mathbb Z})$, it is necessary and sufficient
that there exist $u_0\in\mathrm{Sl}(2,\widehat{\mathbb Z})$ such that
$u^{\unsure{\deleted{0}|\circ}}u_0^{-1}$ commutes with $\mathrm{Sl}(2,\widehat{\mathbb Z})$, and this means
that it is a scalar -- i.e.
$u\in\mathrm{Sl}(2,\widehat{\mathbb Z})\cdot\widehat{\mathbb Z}^{*}$\unsure{,|.}
In other words
the homomorphism (5) is \uline{injective}, and consequently,
\margin{22}

\scanpage{24}
\mspage{300}

the kernel of (4) is none other than the
kernel of the composite
\[
\tag{6}
\Gamma_{\mathbb Q}\xrightarrow{\chi}\widehat{\mathbb Z}^{*}
\longrightarrow\widehat{\mathbb Z}^{*}/\widehat{\mathbb Z}^{*2}
\]
i.e. (by class field theory for
$\mathbb Q$) the \inserted{outer} operation of $\Gamma_{\mathbb Q}$ on
$\mathrm{Sl}(\unsure{1|2},\widehat{\mathbb Z})$ factors
through an \inserted{outer} operation of
$(\Gamma_{\mathbb Q\,\mathrm{ab}})_2$ -- the
``multiquadratic'' \unsure{quotient} of $\Gamma_{\mathbb Q}$ -- and the
latter is faithful.

\margin{\uline{N.B.} The outer operation of $\Gamma_{\mathbb Q}$ on
$\mathrm{Sl}(2,\mathbb Z/\unsure{n|m}\mathbb Z)$ takes place through
the \unsure{intermediary} of
$\unsure{((\mathbb Z/\unsure{n|m}\mathbb Z)^{*})_2|(\mathbb Z/\unsure{n|m}\mathbb Z)_2^{*}}$}

Thus, if we identify $\Gamma_{\mathbb Q}$ \deleted{\unsure{by (4)}}
with a subgroup of $\mathrm{Autext}\,\mathrm{Sl}(2,\widehat{\mathbb Z})^{\wedge}$,
this subgroup \inserted{$\unsure{\Gamma}$} is subject to
fairly strong conditions : \deleted{\illegible{1}} leave invariant
the normal subgroup kernel of
\[
\tag{7}
\mathrm{Sl}(2,\mathbb Z)^{\wedge}\longrightarrow
\mathrm{Sl}(2,\widehat{\mathbb Z})
\]
and induce on $\mathrm{Sl}(2,\widehat{\mathbb Z})$
\deleted{\unsure{the}}\inserted{an} outer operation
coming from a homomorphism $\Gamma\longrightarrow
\widehat{\mathbb Z}^{*}/\widehat{\mathbb Z}^{*2}$.
We can moreover make this homomorphism precise
as follows. Let us consider in $\mathrm{Sl}(2,\widehat{\mathbb Z})$ the
``loop-subgroup'' coming from the point at
\margin{23}

\scanpage{25}

infinity of $M_{1,1\,\mathbb Q}$, it corresponds
to a class of outer homomorphisms
\[
(8)\qquad
T=T_{\infty}(\overline{\mathbb Q}_{0})
\longrightarrow
\mathrm{Sl}(2,\overset{\unsure{\deleted{\wedge}|\wedge}}{\mathbb Z})^{\wedge}
\]
whose image \inserted{in $\mathrm{Sl}(2,\widehat{\mathbb Z})$}
unless I am mistaken, is (up to conjugation) a ``unipotent''
subgroup
\[
(9)\qquad
\left(\begin{array}{cc}1&0\\ n&1\end{array}\right),
\qquad n\in\widehat{\mathbb Z}.
\]

This outer homomorphism is respected by the operation of $\Gamma$, whence a
well-determined operation of $\Gamma$ on \deleted{$\unsure{\mathrm{Sl}(2,}$} $T$
i.e. a homomorphism
\deleted{$\Gamma\longrightarrow\operatorname{Aut}(T)\simeq\widehat{\mathbb Z}^{*}$}
\deleted{(which \illegible{1} will be, precisely $\unsure{\chi}$, the}
\[
(10)\qquad
\Gamma\xrightarrow{\chi}\operatorname{Aut}(T)
\simeq\widehat{\mathbb Z}^{*},
\]
which will be precisely the cyclotomic character.
More precisely, for \unsure{an|An} outer automorphism \inserted{$u$} of
$\mathrm{Sl}(2,\mathbb Z)^{\wedge}$ (or only of $\mathrm{Sl}(2,\widehat{\mathbb Z})$)
which leaves invariant \unsure{the|the} \deleted{\unsure{image}} \unsure{conjugacy|conjugacy} \unsure{class}
of the image of (8) [or \unsure{only|only} of the composite
$T\unsure{\longrightarrow}\mathrm{Sl}(2,\widehat{\mathbb Z})$],
\inserted{we check that} there exists a unique
$\chi(u)\in\operatorname{Aut}(T)\simeq\widehat{\mathbb Z}^{*}$,
such that we have commutativity in
\margin{do the verification}
\margin{24}

\scanpage{26}
\mspage{301}

\[
(11)\qquad
\diagram{\begin{tikzcd}[ampersand replacement=\&]
T \arrow[r, "\unsure{\kappa|k}"] \arrow[d, "\chi(u)"]
 \& \mathrm{Sl}(2,\mathbb Z)^{\wedge} \arrow[d, "u"] \\
T \arrow[r]
 \& \mathrm{Sl}(2,\mathbb Z)^{\wedge}
\end{tikzcd}}
\]
--- or \unsure{even}, only in
\[
(12)\qquad
\diagram{\begin{tikzcd}[ampersand replacement=\&]
T \arrow[r, "\unsure{\kappa|k}"] \arrow[d, "\chi(u)"]
 \& \mathrm{Sl}(2,\widehat{\mathbb Z}) \arrow[d, "\dot u"] \\
T \arrow[r]
 \& \mathrm{Sl}(2,\widehat{\mathbb Z})
\end{tikzcd}}
\quad ---
\]
whence a character \inserted{$\chi$} (10) on the
group of all \inserted{outer} automorphisms $u$ of
$\mathrm{Sl}(2,\widehat{\mathbb Z})$ which have the property \unsure{in question}.

That said, $\Gamma$ is contained in this subgroup \inserted{$\unsure{\Gamma}$},
and moreover in the subgroup of outer automorphisms which \uline{moreover}
preserve the kernel of (7), and which operate \inserted{(outerly!)} in the quotient
$\mathrm{Sl}(2,\widehat{\mathbb Z})$ according to the composite
\[
\Gamma\unsure{:}\xrightarrow{\chi}\widehat{\mathbb Z}^{*}
\longrightarrow\widehat{\mathbb Z}^{*}/\widehat{\mathbb Z}^{*2}
\hookrightarrow\operatorname{Autext}_{\mathrm{lac}!}
(\mathrm{Sl}(2,\widehat{\mathbb Z})).
\]

\deleted{One} Besides, let us note that the center $\{\pm1\}$
of $\mathrm{Sl}(2,\mathbb Z)^{\wedge}$ not being trivial, the outer operation
\margin{25}

\scanpage{27}

of $\Gamma_{\mathbb Q}$ on
$\widehat{\mathcal T}_{1,1}^{+}$
does not allow us a priori to reconstruct (up to unique isomorphism)
the extension $N_{1,1}$ of $\Gamma_{\mathbb Q}$ by
$\widehat{\mathcal T}_{1,1}^{+}$
of (1). It gives only the extension \deleted{\illegible{1}}
$N'_{1,1}$ of $\Gamma_{\mathbb Q}$ by
\[
\widehat{\mathcal T}_{1,1}^{+\prime}
\simeq\mathrm{Sl}(2,\mathbb Z)^{\wedge}/\{\pm1\}
\inserted{\overset{\unsure{\mathrm{dfn}|\mathrm{déf}}}{=}
\mathrm{Sl}(2,\widehat{\mathbb Z})'}
\]
(by inverse image of the extension

\diagram{\begin{tikzcd}[ampersand replacement=\&,column sep=1.1em,row sep=1.0em]
{(13)\quad 1}\arrow[r] \&
\mathrm{Sl}(2,\widehat{\mathbb Z})'\arrow[r] \&
\operatorname{Aut}\deleted{\unsure{\mathrm{ext}}}_{\mathrm{lac}}\bigl(\mathrm{Sl}(2,\unsure{\mathbb Z})^{\wedge}\bigr)\arrow[r] \&
\operatorname{Autext}_{\mathrm{lac}}\,\unsure{\mathrm{Sl}(2,\mathbb Z)^{\wedge}}\arrow[d]\\
\&\&\&1
\end{tikzcd}}

This extension \inserted{$N'_{1,1}$} has also been interpreted
as the extension $M_{0,3}\simeq N'_{0,4}$ of
$\Gamma_{\mathbb Q}$ by
$\widehat{\mathfrak S}_{0,3}
\simeq\widehat{\mathcal T}_{0,4}'$
(group which is itself an extension of
$\mathfrak S_3\simeq\mathrm{Sl}(2,\mathbb Z/2\mathbb Z)$
by $\widehat{\Pi}_{0,3}$).
We already know that the \inserted{outer} operation of $\Gamma_{\mathbb Q}$
on the kernel $\widehat{\Pi}_{0,3}$ of
$\mathrm{Sl}(2,\overset{\unsure{\deleted{\wedge}|\wedge}}{\mathbb Z})^{\wedge\prime}
\longrightarrow\mathrm{Sl}(2,\unsure{\mathbb Z/2|\mathbb Z_{2}})$
is faithful, and commutes with the outer operation of
$\mathfrak S_3\simeq\mathrm{Sl}(2,\mathbb Z/2\mathbb Z)
\simeq\mathrm{Sl}(2,\mathbb Z)^{\wedge\prime}/\widehat{\Pi}_{0,3}$.
But this commutation can at pres-
\margin{26}

\scanpage{28}
\mspage{302}
ent be interpreted as a consequence of the fact that the outer operation of $\Gamma_{\mathbb Q}$ on $\mathfrak S_3\simeq \mathrm{Sl}(2,\mathbb Z/2\mathbb Z)$ is trivial, itself \inserted{(from the present viewpoint)} a consequence of the fact that it \unsure{takes place} through the intermediary of a (faithful) operation of $(\mathbb Z/n\mathbb Z)^*_{\unsure{/2|2}}$, now for $n=2$ $(\mathbb Z/n\mathbb Z)^*_2=1$ (a relation which is on the other hand false for \deleted{\unsure{the}} \inserted{\unsure{other|other}} \unsure{values}, $n\neq\deleted{\unsure{2,}}1,2$ i.e. $n\geq3$). This commutation relation, of a fairly trivial nature, thus appears at present as a very special case of the markedly subtler congruential relations. We can \inserted{try to} express them \deleted{\unsure{in}} \deleted{a way} directly in terms of the outer operation of $\Gamma_{\mathbb Q}$ on $\widehat\Pi_{0,3}$, by introducing the kernels of the \unsure{homomorphisms} $\widehat{\mathcal T}_{1,1}^{\unsure{+\prime|+/n}}\longrightarrow \mathrm{Sl}(2,2m)$, which are open subgroups \deleted{$\Pi(2m)$} :
\[
(14)\qquad \Pi(2m)\subset\widehat\Pi_{0,3}
\]
normal and invariant under $\mathfrak S_3$, kernels
\margin{27}

\scanpage{29}
of homomorphisms
\[
(15)\qquad \varphi_{2m}:\widehat\Pi_{0,3}\longrightarrow\mathrm{Sl}(2,\mathbb Z/2m)'
\]
which it should \unsure{of them|not} be very difficult to make explicit, \deleted{\unsure{in terms}} \inserted{by ``reduction modulo $2m$''} of a \unsure{homomorphism}
\[
(16)\qquad \varphi_{\infty}:\widehat\Pi_{0,3}\longrightarrow\mathrm{Sl}(2,\widehat{\mathbb Z})'
\]
which has already been mentioned previously -- \unsure{homomorphism} whose image is the kernel subgroup of $\mathrm{Sl}(2,\widehat{\mathbb Z})'\longrightarrow\mathrm{Sl}(2,\mathbb Z/2\mathbb Z)$, i.e. we have an exact sequence
\[
(17)\qquad \deleted{1\longrightarrow\Pi(\unsure{2m})\longrightarrow}\inserted{1\longrightarrow\Pi(\infty)\longrightarrow}\widehat\Pi_{0,3}\longrightarrow\mathrm{Sl}(2,\widehat{\mathbb Z})'\longrightarrow\mathrm{Sl}(2,\mathbb Z/2\mathbb Z)\longrightarrow1
\]
\[
(18)\qquad 1\longrightarrow\Pi(2m)\longrightarrow\widehat\Pi_{0,3}\longrightarrow\mathrm{Sl}(2,\mathbb Z/2m\mathbb Z)'\longrightarrow\mathrm{Sl}(2,\mathbb Z/2\mathbb Z)\longrightarrow1.
\]
This being laid down, \unsure{these|the} exact sequences allow us in a natural way to define an outer operation of $\widehat{\mathbb Z}^{*}/\widehat{\mathbb Z}^{*2}$ on $\widehat\Pi_{0,3}/\Pi(\infty)$, or again of the $(\mathbb Z/2m\mathbb Z)^*_2$ on the $\widehat\Pi_{0,3}/\Pi(2m)$, and the ``condition of \unsure{congruence|congruences}'' is expressed by saying that the cyclotomic character $\chi$ of $\Gamma_{\mathbb Q}$ (\inserted{which we can regard as} deduced from the outer operation of
\margin{28}

\scanpage{30}
\mspage{303}
$\Gamma_{\mathbb Q}$ on \unsure{the|the} \unsure{group|groups} with loops $\widehat\Pi_{0,3}$), which gives \unsure{a|in} composite
\[
(19)\qquad \Gamma_{\mathbb Q}\xrightarrow{\chi}\widehat{\mathbb Z}^{*}\xrightarrow{\mathrm{can.}}\widehat{\mathbb Z}^{*}/\widehat{\mathbb Z}^{*2},
\]
\unsure{of|hence} an outer operation of $\Gamma_{\mathbb Q}$ on $\widehat\Pi_{0,3}/\Pi(\infty)$, is \uline{compatible} with the outer operation of $\Gamma_{\mathbb Q}$ on $\widehat\Pi_{0,3}$ itself.

For each integer $m$, this condition is made explicit by the fact that $\Pi(2m)$ is invariant under the outer operation of $\Gamma_{\mathbb Q}$, and that the induced outer operation \unsure{of|on} $\widehat\Pi_{0,3}/\Pi(2m)$ is the one deduced from the character obtained by composing (19) with
\[
\widehat{\mathbb Z}^{*}/\widehat{\mathbb Z}^{*2}\longrightarrow(\mathbb Z/\unsure{2m|m}\mathbb Z)^*_2:
\]
\[
(20)\qquad \Gamma_{\mathbb Q}\longrightarrow(\mathbb Z/m\mathbb Z)^*_2.
\]
We should look whether this formulation is really \uline{equivalent} to that of the conditions of \unsure{congruences|congruence} in $\widehat{\mathcal T}_{1,1}^{+\prime}$ in its \unsure{full} entirety, \deleted{(\unsure{even if it means there} \illegible{1}} \inserted{and not only a} consequence; \unsure{exercise|exercises} ! On the other hand, \unsure{the|the}
\margin{29}
